\head{原著勘误与增补（清单二）}
\source{第三卷第二部分，上传PDF第84--90页；原分册第85--91页，通卷页码217--223}
\NewNote{本章是原作者刊载的勘误全文译文，不是本次核查报告。所有页、行定位均指原版；本合订本页码已改变。原版 A 部分没有连续条号，这里仅按原有条目分段。B 部分保留 $\operatorname{Err}_{\mathrm{III}}$ 的1--30号。涉及难辨字形，均以底本图像而非抽取文本为准。}

\section*{A．排印错误}
\addcontentsline{toc}{section}{原勘误清单二：A．排印错误}
\newcommand{\typo}[2]{\par\smallskip\noindent\textbf{（#1）}\quad #2\par}
\typo{$0_{\mathrm I}$，3.2.1}{第27页第2行，将 $X$ 改为 $U$。}
\typo{$0_{\mathrm I}$，3.2.6}{第27页倒数第1行及第28页第1行，将 $\mathfrak H_\lambda$ 改为 $\mathscr H_\lambda$（共三处）。}
\typo{$0_{\mathrm I}$，3.4.5}{在“3.4.5)”之前加一个左括号。}
\typo{$0_{\mathrm I}$，4.2.1}{第39页第10行，将 $\psi_*(A)$ 改为 $\psi_*(\A)$。}
\typo{$0_{\mathrm I}$，5.1.3}{第45页第16行，将 $\F|V$ 改为 $\F|U$。}
\typo{$0_{\mathrm I}$，5.3.9}{第48页第12行，在“使得”之前加上“在 $x$ 的一个开邻域 $U$ 上”。}
\typo{$0_{\mathrm I}$，7.6.9}{第73页倒数第3行，将 $J_\lambda$ 改为 $\mathfrak J_\lambda$。}
\typo{I，1.1.15}{第83页第3行，将 $\mathfrak a$ 改为 $\mathfrak a\ne A$。}
\typo{I，1.4.1}{第90页倒数第2行，将 $\widetilde N$ 改为 $N$。}
\typo{I，2.3.1}{第101页第3行，将（0，4.1.6）改为（0，4.1.7）。}
\typo{I，2.3.2}{第101页第11行，将 $B$ 改为 $B_s$，将 $C$ 改为 $C_t$。}
\typo{I，2.5.5}{第104页第19行，将“$S$-态射”改为“$S$-预概形”。}
\typo{I，3.3.7}{第109页第17行，将 $Y_{(S)}$ 改为 $Y_{(S')}$。}
\typo{I，3.4.5}{第113页第9行，将 $\underline s$ 改为 $s$。}
\typo{I，3.4.8}{第114页第4行，将 $p(x)$ 改为 $f(x)$。}
\typo{I，3.5.1}{第115页第3行，将 $1_X\times g$ 改为 $1_{X'}\times g$。}
\typo{I，3.7.2}{第119页第8行，将第一个 $X$ 改为 $X'$。}
\typo{I，4.2.2}{第123页第8行，在交换图的左侧竖箭头处，将 $\alpha_{\psi(y)}$ 改为 $\rho_{\psi(y)}$。}
\typo{I，4.4.3}{第126页第16行，将法文误拼“isomorphee”改为“isomorphes”。}
\typo{I，4.5.5}{第127页倒数第17行，将（4.2.4）改为（4.2.5）。}
\typo{I，5.1.4}{第128页倒数第14行，将（2.1.7）改为（2.1.8）。}
\typo{I，5.3.13}{第134页第20行，将（4.2.4）改为（4.2.5）。}
\typo{I，6.4.2}{第147页倒数第1行，将“覆盖”改为“有限覆盖”。}
\typo{I，9.1.13}{第171页第15行，将 $p'^{-1}(\F)$ 改为 $p'^*(\F)$。}
\typo{I，9.5.11}{第179页倒数第7行，将 $Y$ 改为 $Y'$。}
\typo{I，9.6.5}{第180页倒数第17行，将（0，4.1.4）改为（0，4.1.3）。}
\typo{I，10.12.2}{第206页倒数第14行，将 $(X_n)_{(S_n)}$ 改为 $(X_n)_{(S_m)}$。}
\typo{第一卷术语索引}{第219页倒数第6行，将0，4.1.4改为0，4.1.5。第222页倒数第16、17、18行，将I，2.1.7改为I，2.1.8；第222页倒数第19行，将I，2.1.8改为I，2.1.7。}
\typo{II，1.5.2}{第12页第17行，在交换图左侧竖箭头旁加上 $f$。}
\typo{II，4.2.7}{第75页倒数第16行，将（III，2.1.14）改为（III，2.1.13）。}
\typo{II，5.2.1}{第97页倒数第16行，将 $X-U$ 改为 $U$。同页倒数第8行，将 $X^f$ 改为 $X_f$。}
\typo{II，5.3.6}{第100页第16行，将（4.6.18）改为（4.6.17）。}
\typo{II，6.2.7}{第116页倒数第9行，将“第V章”改为“第IV章”。}
\typo{II，6.6.5}{第132页倒数第15行，将“第V章”改为“第IV章”。}
\typo{II，7.4.12}{第151页第1行，将“第V章”改为“第III章”。}
\typo{II，8.1.4}{第154页第15行，将（III，2.3.8）改为（III，2.3.7）。}
\typo{II，8.10.5}{第187页倒数第17行，将 $g_{(j)x}$ 改为 $g_{j(x)}$。}
\typo{$0_{\mathrm{III}}$，10.2.7}{第20页第1行，将“$\mathfrak u$-进制”改为“$\mathfrak n$-进制”；同页倒数第13行，将 $k$ 改为 $k_i$。}
\typo{$0_{\mathrm{III}}$，11.1.1}{第23页第18行，将 $\Im(d_r^{p+r,q-r+1})$ 改为 $\Im(d_r^{p-r,q+r-1})$。}
\typo{$0_{\mathrm{III}}$，11.8.6}{第45页倒数第14行，将箭头 $\to$ 改为 $\rightsquigarrow$（四处）。}
\typo{$0_{\mathrm{III}}$，12.3.4}{第61页第4行，将 $\xrightarrow{\sim}$ 改为 $\to$。}
\typo{$0_{\mathrm{III}}$，13.2.4}{第67页倒数第2行，将 $\varinjlim$ 改为 $\varprojlim$（两处）。}
\typo{III，1.1.7}{第85页第15行，将 $A^r$ 改为 $\Lambda(A^r)$（两处）。}
\typo{III，1.2.4}{第87页第7行，将 $\Gamma(U,\F)$ 改为 $\Gamma(U,\OO_X)$，并将 $\G$ 改为 $\F$；同页第8行，将 $H^p(\mathfrak U,\G)$ 改为 $H^p(\mathfrak U,\F)$。}
\typo{III，2.5.3.1}{第111页第3行，将 $n>0$ 改为 $n\geq0$；同页第8行，将 $-r\leq n\leq0$ 改为 $-r\leq n<0$。}
\typo{III，3.1.2}{第116页第13行，将 $\G$ 改为 $\G\in\mathbf K'$。}
\typo{III，4.3.6}{第133页第7行，将 $K'=K\otimes_A\widehat A$ 改为 $K'\supset K\otimes_A\widehat A$；同页第16行，将 $R(X)$ 改为 $R(Y)$。}
\typo{III，4.4.12}{第138页第10行，将“第V章”改为“第IV章”。}
\typo{III，4.6.6}{第142页第19行，将“第V章”改为“在随后的一节中”。}

\section*{B．正文修改}
\addcontentsline{toc}{section}{原勘误清单二：B．正文修改}
\orig{为便于查找引文，今后将插入第N章的勘误清单中的正文修改，记为 $\operatorname{Err}_{N}$ 后接一个号码。}
\newcommand{\errIII}[1]{\par\addvspace{.6\baselineskip}\Needspace{3\baselineskip}\noindent\hypertarget{errIII:#1}{}\textbf{（$\operatorname{Err}_{\mathrm{III}}$，#1）}\quad\ignorespaces}

\errIII{1}在（$0_{\mathrm I}$，5.1.3）第45页第18行，将“直和”改为“有限直和”。

\errIII{2}在（$0_{\mathrm I}$，5.4.3）第49页倒数第8至第4行，自“在一般情形……”起，替换为下文：

在一般情形，设 $X$ 是赋环空间，每个 $x\in X$ 的 $\OO_x$ 都是局部环；若 $\LL$ 是有限型 $\OO_X$-模，且存在 $\OO_X$-模 $\F$，使 $\LL\otimes_{\OO_X}\F$ 同构于 $\OO_X$，则前面的论证说明，对每个 $x$，$\LL_x$ 都同构于 $\OO_x$。由此可知 $\LL$ 可逆。事实上，取 $x$ 的开邻域 $U$，使 $\LL|U$ 由 $U$ 上的 $n$ 个截面 $s_i$ 生成（5.2.1）。例如可设 $s=s_1$ 满足 $s_x\ne0$。\cmt{C3-A03}{底本这里写的是芽 $s_x\ne0$。仅非零并不足以推出它生成秩一自由模；所需的是选择一个在剩余模中非零的生成元，即 $s(x)\ne0$。这种选择由有限生成元在剩余模中仍生成而得到，因此可补足论证，不是否定本条结论。}由于 $\LL_x\simeq\OO_x$，对每个 $i$，在 $x$ 的某开邻域 $V_i\subset U$ 上存在 $\OO_X$ 的截面 $t_i$，使 $(t_is)_x=(s_i)_x$。因此存在开邻域 $V\subset\bigcap_{i=1}^nV_i$，使 $V$ 上 $t_is=s_i$，即 $\LL|V$ 由单个截面 $s|V$ 生成。再设 $z$ 为由 $s$ 定义的同态 $\OO_X|V\to\LL|V$ 的核在某开集 $W\subset V$ 上的截面（5.1.1）。对每个 $y\in W$，$z_y$ 零化 $\LL_y$，因而由假设 $z_y=0$，于是 $z=0$。断言得证。再考察 $\LL^{-1}\otimes\LL\otimes\F$，立即得到 $\F\simeq\LL^{-1}$。

\errIII{3}在（$0_{\mathrm I}$，7.2.4）中，必须增加条件：各幂 $\mathfrak J^n$ 是可容许环 $A$ 中的闭理想，该命题才正确。已给出的证明不正确；修正后的命题可由 Bourbaki，\textit{Top. gén.}，第III章，第3版，§3，第5目，命题9的推论1推出。在（$0_{\mathrm I}$，7.2.5）与（$0_{\mathrm I}$，7.2.6）的陈述中，同样须假设各 $\mathfrak J^n$ 在 $A$ 中闭。

\errIII{4}在命题（I，2.2.5）之后增加：沿用（I，2.2.5）的记号，若 $\mathfrak J$ 是 $A$ 的理想，则用 $\mathfrak J\OO_X$ 表示（0，4.3.5）中定义的 $\OO_X$-子模。

\errIII{5}在（I，3.4.5）第112页倒数第1行，将“一个域 $K$”改为“一个代数闭域 $K$”。第113页第1行，将“扩张”改为“代数闭扩张”。第113页第1至第4行，从“$K$ 称为……”起，替换为：

对于 $X$ 的任意一个取值于域 $K$ 的点，称 $K$ 为该点的值域；若 $x$ 是这个点的所在地，也说该点位于 $x$。由此定义映射 $X(K)\to X$，将一个取值于 $K$ 的点对应到它的所在地。

第113页第7与第16行删去“几何”；同页倒数第10与第13行也删去“几何”。第115页倒数第10与第11行删去“几何”。

\errIII{6}在（I，3.6.1）证明末尾，第117页倒数第12行，增加：若置 $X'=X\times_Y\Spec(\OO_y/\mathfrak a_y)$，则对每个经 $p$ 与 $X$ 的点相识别的 $x\in X'$，有
\[
 \OO_{X',x}=\OO_{X,x}/\mathfrak a_y\OO_{X,x}.
\]
事实上，问题在 $X,Y$ 上局部，故可设 $X=\Spec B$、$Y=\Spec A$；由平坦性，有 $(B/\mathfrak a B)_x=B_x/\mathfrak a B_x$（0，1.3.2）。

\errIII{7}在（I，5.1.1）第127页倒数第3行，将“拟凝聚 $\OO_X$-模”改为“$\B$ 的拟凝聚理想”。

\errIII{8}在（I，5.1.10）之后，第131页倒数第16行，增加：

\noindent\textbf{注记（5.1.11）。}对（5.1.9）的论证略作调整即可证明：不必预先假设 $X$ 是预概形，结论仍成立；只须假设 $(X,\OO_X)$ 是局部赋环空间，$\I$ 为 $\OO_X$ 的理想，$\I^n=0$，赋环空间 $X_0=(X,\OO_X/\I)$ 为仿射概形，且各 $\I^j/\I^{j+1}$ 为拟凝聚 $\OO_{X_0}$-模。在论证中，只须用（1.8.1）（见 $\operatorname{Err}_{\mathrm{II}}$）代替（2.2.4）。

\errIII{9}在（I，5.3.1）第132页倒数第11行，于“或 $\Delta_X$”之后插入“或 $\Delta_\varphi$，其中 $\varphi:X\to S$ 为结构态射”。

\errIII{10}在（I，5.3.9）中，所给证明不足，因为没有证明 $\Delta_X(X)$ 在 $X\times_SX$ 中局部闭。要得到正确证明，只须使用（4.2.4，a）：对每个 $x\in X$ 及 $x$ 在 $X$ 中的任意仿射邻域 $U$，$U\times_SU$ 是 $\Delta_X(x)$ 的仿射邻域；考虑（5.3.16）（其证明仅使用定义（5.3.1.1）），即可将（5.3.9）归结为 $S=\Spec B$、$X=\Spec A$ 均仿射的情形。此时由（5.3.1.1），$\Delta_X$ 对应于典范同态 $A\otimes_BA\to A$，$x\otimes y\mapsto xy$；这个同态满射，故 $\Delta_X$ 是闭浸入（4.2.3），证明完毕。
\cmt{C3-A02}{上句关于 $U\times_SU$ 仿射，须配合先将底限制到包含 $f(x)$ 的仿射开集，再取映入该开集的仿射 $U$；仅由 $U$ 仿射而不限制 $S$，并不能自动推出任意 $U\times_SU$ 仿射。此处保留作者勘误原文，在核查报告中写出完整局部化步骤。}

\errIII{11}在（I，7.1.15）及（I，7.1.16）第158页倒数第7及第15行，删去“几何”。

\errIII{12}将（I，7.4.7）替换为：将（7.4.1）的定义推广（略有滥用术语）到以下情形：$X$ 是约化预概形，且每个点都有一个仅含有限个不可约分支的开邻域。由（7.3.4）、（7.4.6）可知，对有限型拟凝聚 $\OO_X$-模 $\F$，称其为挠层，等价于 $\Supp(\F)$ 不包含 $X$ 的任何不可约分支。

\errIII{13}在（I，10.11.7）第205页第17至第23行，从“尚须证明……”起的证明结尾无须保留：由（0，5.3.5），所考虑的层都是凝聚层，故可直接应用（10.11.6）。

\errIII{14}在（II，1.7.8）第16页倒数第10行之后增加：当 $\E=\OO_S^n$ 时，也把 $\mathbf V(\E)$ 写成 $\mathbf V_S^n$；若 $S=\Spec A$ 仿射，则把 $\mathbf V_S^n$ 写成 $\mathbf V_A^n$；有 $\mathbf V_A^n=\Spec(A[T_1,\ldots,T_n])$，其中 $T_i$ 为不定元。

\errIII{15}在（II，1.7.10）第17页第10行及第12、16--17行，删去“几何”。

\errIII{16}在（II，1.7.12）中增加：同样约定
\[
 \mathbf V(\OO_S^n)=\mathbf V_S^n=S[T_1,\ldots,T_n].
\]

\errIII{17}在（II，4.2.6）第75页第7、8行，删去“几何”。

\errIII{18}在（II，4.6.13）第92页第7行，须将论证说得更精确：沿用（4.4.10）的记号，这里须证明浸入 $\Gamma_f$ 拟紧，才能应用（i bis）。由（4.6.4），为证明 $\LL$ 关于 $f$ 丰沛，可设 $Y$ 仿射。又在（v）的每一种假设下，$f$ 都拟紧（I，6.6.4）。若 $g$ 分离，则 $\Gamma_f$ 为闭浸入（I，5.4.3），因而拟紧（I，6.6.4）。若 $X$ 局部诺特，则因 $Y$ 仿射且 $f$ 拟紧，$X$ 的底空间拟紧，从而诺特；再次应用（I，6.6.4），可知 $\Gamma_f$ 拟紧。

\errIII{19}（II，5.2.2）的陈述不正确：若不知道对 $Y$ 的开集 $U$，每个拟凝聚 $(\OO_X|f^{-1}(U))$-模都是某个拟凝聚 $\OO_X$-模在 $f^{-1}(U)$ 上的限制，就不能说条件 b)、c)、c') 在 $Y$ 上是局部的。因此，应将第98页第16行的“设 $f:X\to Y$ 为分离拟紧态射”替换为：“设 $X,Y$ 为两个预概形，且 $X$ 是概形，或 $X$ 的底空间局部诺特；设 $f:X\to Y$ 为拟紧态射。”证明中应注意：在 $X$ 为概形的情形，由（I，5.5.5、5.5.8），这些假设蕴含 $f$ 分离，因此两种情形下都可应用（I，9.2.2）。

\errIII{20}在（II，6.2.3）第115页倒数第18行之后增加：若存在 $y=f(x)$ 的仿射开邻域 $V$ 及 $x$ 的仿射开邻域 $U$，使 $f(U)\subset V$，且限制态射 $U\to V$ 拟有限，则称态射 $f:X\to Y$ 在点 $x\in X$ 处拟有限。若它在每个 $x\in X$ 处都拟有限，则称 $f$ 局部拟有限。

\errIII{21}（II，6.4.3）的证明不正确：$K$ 的一个有限型 $A$-子模的元素，不一定在 $A$ 上整。将第121页最后九行替换为下文：因为可设 $E$ 无挠，因而忠实，$E$ 也是忠实 $A[u]$-模（这里 $A$ 嵌入 $E$ 的自同态环）。又 $E$ 是有限型 $A$-模，故 $u$ 在 $A$ 上整（Bourbaki，\textit{Alg. comm.}，第V章，§1，第1目，引理1）。立即推出 $u\otimes1$ 的特征值（在 $K$ 的某个代数闭包中）在 $A$ 上整，因而各 $a_i(u)$ 也在 $A$ 上整。

\errIII{22}在（$0_{\mathrm{III}}$，9.1.1）第12页倒数第18行，删去“开”。

\errIII{23}在（$0_{\mathrm{III}}$，11.7.3）第43页倒数第10行，于“$C''$”之后增加：使得对 $C$ 的每个投射对象 $P$（分别对 $C'$ 的每个投射对象 $P'$），函子 $A'\mapsto T(P,A')$（分别 $A\mapsto T(A,P')$）在 $C'$（分别 $C$）中正合。

\errIII{24}（$0_{\mathrm{III}}$，13.7.7）的陈述不准确，须作如下修改：第78页第6行，删去“以及 $(R^{n+1}u(A_i))_{i\in I}$”；同页第7行，将“$R^nu(A)$ 是有限型 $S$-模”改为“$R^nu(A)$ 与 $R^{n+1}u(A)$ 均为有限型 $S$-模”；同页第12--13行，删去“以及 $p+q=n+1$”。证明中，第78页第23行，删去“以及对于 $n+1$”。

\errIII{25}在（III，1.4.15）第92页倒数第6行，将“有限型”改为“拟紧”。

\errIII{26}在（III，2.2.4）第102页第1行，将“$\F$ 与 $\G$ 的支集在 $Y$ 上固有”改为“$\Im(\F\to\G)$ 的支集在 $Y$ 上固有”。证明中，从“既如此……”开始的第10至第16行，替换为下文：

置 $\G_1=\Im(\F\to\G)=\Ker(\G\to\HH)$，它凝聚。由正合列 $0\to\G_1(n)\to\G(n)\to\HH(n)$ 及 $f_*$ 左正合，对每个 $n$，序列 $0\to f_*(\G_1(n))\to f_*(\G(n))\to f_*(\HH(n))$ 正合。因此，只须证明 $n$ 充分大时，$f_*(\F(n))\to f_*(\G_1(n))\to0$ 正合。换言之，可设 $\HH=0$，且 $\G$ 的支集在 $Y$ 上固有，因而在 $Z$ 中闭（II，5.4.10）。于是 $\G'=i_*(\G)$ 为凝聚 $\OO_Z$-模，且 $\G'|X=\G$。由（I，9.4.3），存在凝聚 $\OO_Z$-模 $\F'$，使 $\F'|X=\F$。又因 $X$ 在 $Z$ 中开，且 $\Supp(\G)\subset X$ 在 $Z$ 中闭，可定义满同态 $u':\F'\to\G'$：当开集 $U\subset Z$ 不交 $\Supp(\G)$ 时取 $u'|U=0$，当 $U\subset X$ 时取 $u'|U=u|U$。所得限制 $u'|X$ 正是给定的满同态 $u:\F\to\G$。

既如此，如（2.2.2）证明开头，可设 $Y$ 仿射。只须证明 $n$ 充分大时，$\Gamma(u):\Gamma(X,\F(n))\to\Gamma(X,\G(n))$ 满射。有交换图
\[
\begin{tikzcd}[column sep=large]
\Gamma(Z,\F'(n))\arrow[r,"\Gamma(u')"]\arrow[d,"v"']&\Gamma(Z,\G'(n))\arrow[d,"w"]\\
\Gamma(X,\F(n))\arrow[r,"\Gamma(u)"']&\Gamma(X,\G(n)).
\end{tikzcd}
\]
其中 $v,w$ 是限制同态。由 $\G'$ 的定义，$w$ 双射；由（2.2.3），$n$ 充分大时 $\Gamma(u')$ 满射，故 $\Gamma(u)$ 也满射。

\errIII{27}在（III，2.2.5）第102页倒数第7行之后增加：

（iii）在（2.2.4）关于 $X,Y,f,\LL$ 的假设下，若 $\F$ 是其支集在 $Y$ 上固有的凝聚 $\OO_X$-模，类似（2.2.4）的论证给出整数 $N$，使 $n\geq N$ 时，对所有 $i>0$ 均有 $R^if_*(\F(n))=0$。由上同调正合列可知：若 $u:\F\to\G$ 为凝聚 $\OO_X$-模的同态，且 $\Ker u$ 与 $\Coker u$ 的支集均在 $Y$ 上固有，则存在 $N$，使 $n\geq N$ 时，对所有 $i>0$，相应同态 $R^if_*(\F(n))\to R^if_*(\G(n))$ 均双射。

\errIII{28}在（III，4.4.9）第137页倒数第17行，将“设 $Y$ 为整的局部诺特预概形”改为“设 $X,Y$ 为两个整的局部诺特预概形”；同页倒数第16行，将“有限型”改为“局部有限型”。证明实质上不变：$f^{-1}(y)$ 在 $k(y)$ 上局部有限型，因而仍是离散的。

\errIII{29}在（I，9.3.4）第173页倒数第19行，将“诺特”改为“拟紧”；同页倒数第18、19行，将“凝聚”改为“有限型拟凝聚”。证明中，归结为 $X=\Spec A$、$\F=\widetilde M$，其中 $M$ 是有限型 $A$-模，$\J=\widetilde{\mathfrak J}$，其中 $\mathfrak J$ 是 $A$ 的有限型理想；其余论证不变。

\errIII{30}将（I，9.3.5）第173页最后七行替换为下文：

\noindent\textbf{命题（9.3.5）。}设 $X$ 为预概形，$\F$ 为有限型拟凝聚 $\OO_X$-模。则存在 $X$ 的闭子预概形 $Y$，其底空间等于 $\Supp(\F)$；还存在有限型拟凝聚 $\OO_Y$-模 $\G$，使得若 $j:Y\to X$ 为典范包含，则 $\F\simeq j_*(\G)$。

只须证明 $\F$ 的零化理想层 $\J$ 拟凝聚。取 $Y$ 为由 $\J$ 定义的闭子预概形（I，4.1.2）；因 $\J\F=0$，$\F$ 是 $\OO_X/\J$-模，取 $\G=j^*(\F)$ 即可。为证明 $\J$ 拟凝聚，可局部化为 $X=\Spec A$、$\F=\widetilde M$，其中 $M$ 由有限个元素 $x_i$（$1\leq i\leq r$）生成。于是 $\J$ 是各 $x_i$ 的零化理想层的交。$x_i$ 的零化理想层是同态 $\OO_X\to\F$ 的核，该同态对应于 $A\to M$、$s\mapsto sx_i$；故它拟凝聚（I，4.1.1），这类理想层的有限交也拟凝聚（I，1.3.10）。
